% Bewertungspaket zum Gesamttest «Funktionstypen und Graphen» — Quelle des PDFs (python3 scripts/build-lp-pdf.py funktionstypen).
% Ganze Punkte je Teilschritt; Werte und Fehlerwerte mit python3 nachgerechnet (scripts/lp/funktionstypen/zahlen.py).
% Gerüst aus dem Bewertungspaket von LP 3 (gleichungen-formeln): Auftrag an die KI, Folgefehler, (E)-Zeilen.
\documentclass[11pt]{article}
\usepackage{../lp-druck}
\usepackage{ft-gitter}
\renewcommand{\lpname}{Funktionstypen und Graphen}
\renewcommand{\lpteilgebiet}{Funktionen ·}
\renewcommand{\lpfuss}{Bewertungspaket}
\newcolumntype{P}{>{\centering\arraybackslash}p{10mm}}
\newenvironment{raster}{\par\smallskip\noindent\renewcommand{\arraystretch}{1.6}\linespread{1.35}\selectfont\tabularx{\linewidth}{@{}>{\raggedright\arraybackslash}p{38mm}P X@{}}\toprule
  \small\textsc{Teilschritt} & \small\textsc{P} & \small\textsc{Musterlösung}\\\midrule}{\bottomrule\endtabularx\par}
\newcommand{\fehler}[1]{\par\smallskip{\small\textcolor{lprot}{\bfseries Typische Fehler:} #1}}
\newcommand{\E}{\,{\footnotesize\bfseries\color{lpgruen}(E)}}
\begin{document}
\lpkopf{Bewertungspaket zum Gesamttest}{}

\begin{lpachtung}{Erst lösen, dann öffnen}
Dieses Paket enthält die Musterlösung. Wer es vor dem Lösen liest, misst am Ende nur, wie gut das Abschreiben gelingt.
\end{lpachtung}

\section*{1 · So gehst du vor}
\begin{enumerate}[leftmargin=*,itemsep=2pt]
\item \textbf{Gesamttest lösen} — vollständig, mit Rechenweg, ohne dieses Paket und ohne Taschenrechner, mit Lineal.
\item \textbf{Selbst bewerten} mit Abschnitt~3: Zeile für Zeile vergleichen und die Punkte eintragen. Das geht ohne Hilfsmittel
  und ist der Normalfall.
\item \textbf{Oder bewerten lassen.} Ob du dafür eine KI nutzt und welche, entscheidest du selbst und in eigener Verantwortung —
  je nachdem, was dir aufgrund deines Alters und deiner persönlichen Situation erlaubt ist (Altersgrenzen und
  Nutzungsbedingungen des Anbieters, allenfalls Einverständnis deiner Eltern). Fotografiere oder scanne deine Lösung gut lesbar,
  \textbf{ohne Namen}. Handyfotos tragen oft unsichtbar Standort, Zeit und Gerät mit (Metadaten) — lade lieber einen Scan oder ein
  Bildschirmfoto des Fotos hoch. Lade nur deine Lösung, das Testblatt und dieses PDF hoch, keine persönlichen Daten, und füge den Auftrag aus
  Abschnitt~2 ein.
\item \textbf{Rückmeldung prüfen.} Stimmt, was die KI aus deiner Handschrift gelesen hat? Eine KI kann sich irren —
  im Zweifel gilt das Raster in Abschnitt~3.
\item \textbf{Gezielt wiederholen} nach Abschnitt~4.
\end{enumerate}

\section*{2 · Auftrag an die KI}
\begin{tcolorbox}[colback=lplinie!25,colframe=lplinie,boxrule=0.5pt,arc=1mm,fontupper=\small]
Du bist Korrektorin oder Korrektor für Mathematik auf der Sekundarstufe I in der Schweiz. Im Anhang findest du (1)~das Testblatt mit den Aufgaben, (2)~das Bewertungspaket zum Gesamttest «Funktionstypen und Graphen» mit Musterlösung und Punkteraster und (3)~Bilder meiner handschriftlichen Lösung.\par\medskip
Geh so vor:\par
1. Schreib zu jeder Aufgabe G1 bis G8 zuerst ab, was du in meiner Lösung liest — Rechenweg, Ergebnis und Begründung; bei Skizzen (G2, G8), wie sie verlaufen. Was du nicht sicher lesen kannst, markierst du mit [unsicher]. Nicht raten. Achte besonders auf Minuszeichen, Dezimalpunkte und Einheiten.\par
2. Vergib die Punkte genau nach dem Raster im Bewertungspaket (Abschnitt 3), ganze Punkte je Zeile. Ziehe einen Fehler nur einmal ab: Rechne ich mit einem falschen Zwischenergebnis richtig weiter, gibt es die Folgepunkte. Ausnahme: Zeilen mit \textbf{(E)} gibt es nur für das richtige Ergebnis, nie als Folgepunkt. Die Zeilen «Typische Fehler» sagen, welche Rasterzeile bei diesem Fehler 0 Punkte gibt — sie sind kein zusätzlicher Abzug.\par
3. Andere richtige Begründungen zählen voll, zum Beispiel ein anderer unterscheidender Testpunkt oder ein anderes Merkmal des Graphen. Antworten in Worten zählen, wenn sie dasselbe sagen wie die Musterlösung. Gleichwertig: $y = \tfrac{72}{x}$ und $y = 72 : x$; $d = 12 - 4t$ und $d = -4t + 12$; Brüche und Dezimalzahlen mit gleichem Wert; m/min und eine richtig umgerechnete Einheit. Eine Zuordnung ohne Begründung gibt keinen Punkt.\par
4. Begründe jeden Abzug in einem Satz und nenne die Fehlerart (zum Beispiel «Parabel und Gerade verwechselt», «proportional statt linear mit Anfangswert», «Stand statt Steigung gelesen», «zurückgelegt und Entfernung vertauscht»). Schreib die Musterlösung nicht ab — zeig mir, wo mein Fehler liegt.\par
5. Gib zum Schluss eine Tabelle «Aufgabe | Punkte | Begründung», die Gesamtpunktzahl (von 25) und — nach der Selbsteinschätzung in Abschnitt 4 — welches Kapitel des Leitprogramms ich wiederholen soll. Keine Note.\par
6. Fehlt ein Bild oder ist eine Aufgabe unlesbar, sag es, statt sie zu bewerten.
\end{tcolorbox}

\needspace{16\baselineskip}
\section*{3 · Musterlösung und Punkteraster}
\textbf{Allgemein:} Ganze Punkte je Zeile. Ein Fehler wird nur einmal abgezogen (Folgefehler). Gleichwertige Wege zählen voll.
In der Spalte P: \textbf{(E)}~= Ergebnispunkt, Ergebnis richtig \emph{und} Weg erkennbar, direkt aus den Angaben; ohne Zeichen = die
Zeile bewertet den Weg, die Skizze oder die Begründung. Der ganze Test ist ohne Taschenrechner gestellt.
«Typische Fehler» nennt, welche Zeile 0 Punkte gibt, und die Punkte, die bleiben.
\textbf{Teilrichtige Zeilen:} Verlangt eine Zeile mehrere Angaben und stimmt nur ein Teil, gibt die Zeile 0 Punkte,
ausser die Zeile nennt ausdrücklich eine Aufteilung.

\begin{lpaufgabe}{G1}{Zuordnen}{4 P}
\begin{raster}
A & 1 & $y = 2\sqrt{x}$: beginnt bei $(0 \mid 0)$, nur $x \ge 0$, wird flacher; Testpunkt $x = 4$: $y = 4$\\
B & 1 & $y = 0.5x^2$: Parabel, symmetrisch zur $y$-Achse; Testpunkt $x = 4$: $y = 8$\\
C & 1 & $y = \tfrac{4}{x}$: zwei Äste, trifft keine Achse; Testpunkt $x = 4$: $y = 1$\\
D & 1 & $y = 2x + 1$: Gerade, schneidet die $y$-Achse bei $1$\\
\end{raster}
\fehler{Ein Testpunkt zählt nur, wenn er die gewählte Gleichung von den anderen unterscheidet: $(2 \mid 2)$ liegt auf B \emph{und} C, $(0 \mid 0)$ auf A \emph{und} B. Eine Zuordnung ohne Merkmal oder unterscheidenden Testpunkt: diese Zeile 0. A und B vertauscht (beide durch $(0 \mid 0)$): beide Zeilen 0 $\to$ 2 von 4\,P.}
\end{lpaufgabe}

\begin{lpaufgabe}{G2}{Wertetabelle}{3 P}
\begin{raster}
Tabelle & 1\E & $x = 1,\ 2,\ 4,\ 8$: $y = 8,\ 4,\ 2,\ 1$\\
Skizze & 1 & fallende Kurve durch diese Punkte, trifft die Achsen nicht (nähert sich ihnen)\\
$x = 0$ & 1 & $\tfrac{8}{0}$ gibt es nicht: Durch null kann man nicht teilen.\\
\end{raster}
\fehler{$x \cdot 8$ gerechnet ($8,\ 16,\ 32,\ 64$): Tabellenzeile 0; Skizze als Folgepunkt nur, wenn sie zur eigenen Tabelle passt $\to$ 2 von 3\,P.
Skizze als Gerade oder mit Schnitt einer Achse: Skizzenzeile 0 $\to$ 2 von 3\,P.}
\end{lpaufgabe}

\begin{lpaufgabe}{G3}{Zusammenhänge}{3 P}
\begin{raster}
(a) & 1 & proportional: doppelte Menge, doppelter Preis; $0$ kg kosten nichts ($y = 2.4x$)\\
(b) & 1 & antiproportional: doppelt so schnell, halb so lange; $x \cdot y = 90$\\
(c) & 1 & linear mit Anfangswert: $20$ Fr. auch ohne Telefonieren ($y = 0.1x + 20$)\\
\end{raster}
\fehler{Antwort ohne Begründung: diese Zeile 0 $\to$ 2 von 3\,P. (c) «proportional»: Zeile (c) 0 $\to$ 2 von 3\,P.}
\end{lpaufgabe}

\begin{lpaufgabe}{G4}{Antiproportional}{3 P}
\begin{raster}
(a) & 1\E & $8 \cdot 9 = 72$ Personentage, $72 : 6 = 12$ Tage\\
(b) & 1 & $y = \dfrac{72}{x}$\\
(c) & 1 & $\tfrac{72}{12} = 6$ Tage. Auch als Folgepunkt mit der eigenen Gleichung aus (b).\\
\end{raster}
\fehler{Proportional gerechnet ($6$ Personen: $\tfrac{9}{8} \cdot 6 = 6.75$ Tage, Gleichung $y = \tfrac{9}{8}x$): Zeilen (a) und (b) 0, (c) als Folgepunkt $\to$ 1 von 3\,P.}
\end{lpaufgabe}

\begin{lpaufgabe}{G5}{Velofahrt}{4 P}
\begin{raster}
(a) & 1\E & $\tfrac{480}{6} = 80$ m/min\\
(b) & 1 & Pause: Die Entfernung bleibt $480$ m, der Graph ist waagrecht.\\
(c) & 1\E & $\tfrac{1280 - 480}{8} = 100$ m/min\\
(d) & 1 & von $12$ bis $20$ min: Dort ist der Graph am steilsten. Auch als Folgepunkt passend zu den eigenen Werten aus (a) und (c).\\
\end{raster}
\fehler{Stand statt Steigung ($480$ bzw.\ $1280$ angegeben): diese Zeile 0 $\to$ 3 von 4\,P.
(c) $1280 : 8 = 160$ (Anfangsstand $480$ vergessen): Zeile (c) 0, (d) als Folgepunkt $\to$ 3 von 4\,P.}
\end{lpaufgabe}

\begin{lpaufgabe}{G6}{Welcher Graph?}{2 P}
\begin{raster}
Wahl & 1\E & A\\
Begründung & 1 & B fällt steiler als er gestiegen ist: zurück schneller, nicht langsamer. C steigt nach der Pause weiter: Tim fährt nicht nach Hause. (Beide Gründe verlangt.)\\
\end{raster}
\fehler{Wahl A ohne Begründung zu B und C: Begründungszeile 0 $\to$ 1 von 2\,P.}
\end{lpaufgabe}

\begin{lpaufgabe}{G7}{Wanderung}{3 P}
\begin{raster}
(a) & 1 & $s = 4t$ und $d = 12 - 4t$ (für $0 \le t \le 3$; die Angabe des Bereichs ist nicht verlangt)\\
(b) & 1 & $s = 8$ km, $d = 4$ km. Auch als Folgepunkt mit den eigenen Gleichungen.\\
(c) & 1\E & $12 - 4t = 3$, $4t = 9$, $t = 2.25$ h ($2$ h $15$ min)\\
\end{raster}
\fehler{$s$ und $d$ vertauscht: Zeile (a) 0, (b) 1 als Folgepunkt (mit den eigenen Gleichungen richtig gerechnet), (c) 0 $\to$ 1 von 3\,P.
(c) $3 : 4 = 0.75$ h (zurückgelegt statt Entfernung): Zeile (c) 0 $\to$ 2 von 3\,P. $2.25$ h als «$2$ h $25$ min» umgerechnet: Zeile (c) 0 (richtig: $0.25 \cdot 60 = 15$ min).}
\end{lpaufgabe}

\begin{lpaufgabe}{G8}{Diagramm prüfen}{3 P}
\begin{center}
\wegzeit{6}{4}{5}{200}{zurückgelegt [m]}{\draw[lpblau,line width=1.2pt] (0,0) -- (1,2) -- (2,2) -- (4,3.5);}
\end{center}
\begin{raster}
Fehler 1 & 1 & Senkrechtes Stück bei $t = 5$ min: $300$ m Weg ohne Zeit, also unendlich schnell. Unmöglich.\\
Fehler 2 & 1 & Fallendes Stück: Der zurückgelegte Weg kann nicht kleiner werden, auch auf einem Rückweg nicht.\\
Skizze & 1 & Ein mögliches Diagramm mit Pause: steigend, waagrecht, wieder steigend (wie im Bild); nie fallend, nie senkrecht.\\
\end{raster}
\fehler{«Das fallende Stück heisst, sie wird langsamer»: Zeile Fehler 2 0 $\to$ 2 von 3\,P. Nur einen Fehler gefunden: eine Fehlerzeile 0 $\to$ 2 von 3\,P.}
\end{lpaufgabe}

\needspace{12\baselineskip}
\section*{4 · Selbsteinschätzung}
\begin{tabularx}{\linewidth}{@{}l X@{}}
\toprule
\textbf{22 – 25 P} & Die geprüften Teile sitzen. Wo du Punkte verloren hast: das Kapitel dieser Aufgabe nochmals.\\
\textbf{17 – 21 P} & Den schwächsten Teil nochmals, gemessen am Anteil (Teil A 7, B 6, C 6, D 6 P): Tüfteln und Übungen des Kapitels dieses Teils.\\
\textbf{11 – 16 P} & Zurück zu den Kapiteln aller Aufgaben, in denen du Punkte verloren hast.\\
\textbf{0 – 10 P} & Zurück zu den Kapiteln jedes Teils, in dem du weniger als die Hälfte erreicht hast; sind es alle vier, zum Vortest und zu Kapitel 1.\\
\bottomrule
\end{tabularx}

\smallskip
\textbf{Aufgabe $\to$ Kapitel:} G1, G2 $\to$ 1; G3, G4 $\to$ 2; G5, G6 $\to$ 3; G7, G8 $\to$ 4.
\end{document}
